% !TEX root = ../../../tree_notes.tex
% =====================================================================
%  15.1  The two-DFS template and the edge-transfer rule  sec:c20-twodfs
%  Exports: def:c20-pair-state  thm:c20-edge-transfer  lst:c20-twodfs
%           fig:c20-updown  ex:c20-twodfs-run  prob:c20-twodfs-solved
%           lst:c20-1187e-soln
%  THE parent-side piece Ch.13 promised and Ch.14 refused to maintain.
%  The down pass IS Ch.13's fold (thm:c10-subtree-template), reused; the
%  NEW object is the up pass and the edge-transfer rule that manufactures
%  it.  Anchor problem CF 1187E (Tree Painting): maximise the sum of
%  subtree sizes over every root, F(c)=F(v)+n-2*size(c).  All overlays
%  hand-computed on the running tree T from the opener.
% =====================================================================
\section{The Two-DFS Template: A Down Fold and an Up Fold}\label{sec:c20-twodfs}

\begin{previouslybox}
\begin{enumerate}[nosep, leftmargin=1.6em]
  \item Chapter~13 folded a \emph{constant} tuple up from the children at a
        \emph{fixed} root (\thmref{c10-subtree-template}) and read one number
        at that root. What did it do about the other $n-1$ roots? (nothing --
        it fixed one root and stopped; the parent-side piece was punted here)
  \item Chapter~12 computed the subtree size $\size(v)$ of every node in one
        measuring DFS (\coderef{c09-dfs3}). On $T$ rooted at~$1$, name
        $\size(2)$ and $\size(9)$. ($4$ and $2$)
  \item What is the identity element of a $\bigl(\max\bigr)$-fold? of a
        sum? of an XOR? (of $\max$ there is none unless you adjoin
        $-\infty$; $0$ for sum; $0$ for XOR) -- you will need it to seed the
        up pass.
  \item State Ch.~13's leaf convention again. ($\size(\nil)=0$, so a leaf's
        fold is just its own seed and needs no \code{if})
\end{enumerate}
\end{previouslybox}

\begin{topicintro}
Every problem in this chapter shares one sentence in its output: \emph{a value
per vertex, as if the tree were rooted at that vertex.} Running Ch.~13's DFS
once per vertex costs $\bigO{n^2}$ and dies at $n=2\cdot10^{5}$. The fix is a
single observation: when the root slides from a node $v$ to a neighbour $c$,
\emph{only the edge $(v,c)$ changes which side of it each node sits on}, so the
answer moves by a small, local amount. Turn that local amount into a rule and
two DFS passes -- one bottom-up, one top-down -- produce all $n$ answers in
$\bigO{n}$. The bottom-up pass is Ch.~13's fold verbatim; the top-down pass is
the parent-side piece that chapter~13 announced and never delivered. This
section builds the machinery once; sections 15.2--15.4 only change what the
fold and the shift are.
\end{topicintro}

\subsection{The Down/Up Split}

\begin{definition}{The (down, up) pair state at a node}{c20-pair-state}
Fix a reference root $r_0$ (we use $r_0=1$). For every node $v$ split its
rooted answer into two independent pieces:
\begin{itemize}[nosep,leftmargin=1.4em]
  \item $\mathit{down}[v]$ -- the fold's value contributed by
        $\subtree(v)$, i.e.\ what Ch.~13's \thmref{c10-subtree-template}
        already computes: $\mathit{down}[v] = \Phi\bigl(\text{self}(v),
        \mathit{down}[c_1], \dots\bigr)$.
  \item $\mathit{up}[v]$ -- the fold's value contributed by everything
        \emph{outside} $\subtree(v)$, i.e.\ the part of the tree you reach
        only by going \emph{through $v$'s parent}. For the reference root,
        $\mathit{up}[r_0]$ is $\Phi$'s identity (there is nothing above it).
\end{itemize}
The answer with $v$ itself as root is their join
$\mathit{dp}[v] = \Phi\bigl(\mathit{down}[v],\, \mathit{up}[v]\bigr)$.
\end{definition}

The word \term{down} is geometric: $\mathit{down}[v]$ is computed by a DFS that
visits children before parents (post-order). The word \term{up} is equally
geometric: $\mathit{up}[v]$ is computed by a DFS that visits parents before
children (pre-order), because a node's parent-side value is assembled from its
own parent's answer -- so it can only be known once the parent's is. Two DFS
passes, opposite directions, one tree.

\subsection{The Edge-Transfer Rule}

\begin{theorem}{Moving the root across one edge updates the fold locally}{c20-edge-transfer}
Let $v$ be a node and $c$ one of its children in the reference rooting. Then
\[
  \mathit{up}[c] \;=\; \Phi\text{-}\!\operatorname{shift}\!\Bigl(
     \Phi\bigl(\text{self}(v),\; \mathit{up}[v],\;
     \textstyle\bigsqcup_{c' \in \mathit{ch}(v),\, c' \ne c} \mathit{down}[c']\bigr)
  \Bigr),
\]
i.e.\ $c$'s parent-side value is the fold over \emph{everything that is not in
$\subtree(c)$} -- $v$'s own contribution, $v$'s parent-side value, and all of
$c$'s siblings -- followed by whatever per-edge adjustment $\Phi$-shift the fold
demands. Concretely, if the quantity at the root is $A(v)=\Phi(\mathit{down}[v],
\mathit{up}[v])$ and the fold's combine is a plain sum of per-node sizes, the
rule collapses to the closed form $\mathit{up}[c]=\mathit{up}[v]+(n-\size(c))$.
\end{theorem}

\begin{proofbox}[Only the edge $(v,c)$ re-partitions the tree]
With root $r_0$, node $c$'s answer splits as $\mathit{down}[c]$ (inside
$\subtree(c)$) plus $\mathit{up}[c]$ (outside it). When we instead ask for
$c$ as the root, the set of nodes is partitioned by the single edge $(v,c)$ into
$\subtree(c)$ (unchanged membership and unchanged internal parent--child
orientation) and its complement $V\setminus\subtree(c)$ (also unchanged
internally). Hence the value $c$ accumulates \emph{from its own subtree} is
exactly $\mathit{down}[c]$, and the value it accumulates \emph{from outside}
is the fold over that complement.

But the complement, seen from $v$, is precisely $v$'s own contribution,
$v$'s outside piece $\mathit{up}[v]$, and every child of $v$ other than $c$
(each contributing its $\mathit{down}$). That is the inner $\Phi$ in the
statement. Re-rooting at $c$ additionally flips the direction of the edge
$(v,c)$: $c$ now sits one level closer to every node in $V\setminus\subtree(c)$
and one level further from every node in $\subtree(c)\setminus\{c\}$ -- this
uniform re-indexing is what $\Phi$-shift encodes, and for size/depth-sum folds
it is exactly the $\pm(n-2\size(c))$ correction. Since $\mathit{up}[c]$ depends
only on $\mathit{up}[v]$, $\mathit{down}[\cdot]$ of $c$'s siblings, and
self-values -- all known before we descend to $c$ in the pre-order pass --
processing $v$ before $c$ is well-founded. Induction along the pre-order DFS
proves every $\mathit{up}[\cdot]$. \qed
\end{proofbox}

\begin{insightbox}[Why the direction of the second DFS is forced]
The rule writes $\mathit{up}[c]$ in terms of $\mathit{up}[v]$ and the
$\mathit{down}$ values of $c$'s \emph{siblings}. The siblings' $\mathit{down}$
are finished by the post-order pass, and $\mathit{up}[v]$ is finished earlier in
the same top-down pass. So a node's $\mathit{up}$ becomes available exactly when
we arrive at it from its parent -- which is why the up pass must be
$\text{parent}\to\text{child}$, the mirror of the down pass. Get the direction
backwards and you read $\mathit{up}[v]$ before it exists.
\end{insightbox}

\subsection{The Algorithm: One DFS Down, One DFS Up}

\begin{lstlisting}[caption={The generic two-DFS rerooting template. The down pass is Ch.13's fold verbatim; the up pass applies thm:c20-edge-transfer parent-before-child.}, label={lst:c20-twodfs}]
// Two-DFS rerooting.  down[v] = fold over subtree(v) (Ch.13 fold, verbatim).
// up[v]    = fold over everything OUTSIDE subtree(v), reached via parent.
// ans[v]   = Phi(down[v], up[v])  = "the answer with v as root".
// SEED: up[root] = identity of Phi  (empty combine).  NOT a hard zero unless
//       Phi's identity IS zero -- forgetting this mis-defines every up value.
vector<vector<int>> adj;
vector<Val> down_, up_, ans_;

Val combine(Val a, Val b);                       // Phi: associative fold
Val shift(Val parentAggregate, int v, int c);    // per-edge adjustment

void dfs_down(int v, int p) {                    // post-order (children first)
    down_[v] = self_value(v);                    // v's own seed
    for (int c : adj[v]) if (c != p) {
        dfs_down(c, v);
        down_[v] = combine(down_[v], down_[c]);  // fold one child at a time
    }
}
void dfs_up(int v, int p) {                      // pre-order (parent first)
    ans_[v] = combine(down_[v], up_[v]);         // v's every-root answer
    for (int c : adj[v]) if (c != p) {
        // everything NOT in subtree(c) = up[v], v's self, and c's siblings
        Val rest = up_[v];
        rest = combine(rest, self_value(v));
        for (int s : adj[v])
            if (s != p && s != c) rest = combine(rest, down_[s]);
        up_[c] = shift(rest, v, c);              // thm:c20-edge-transfer
        dfs_up(c, v);
    }
}
// call: dfs_down(1, 0); up_[1] = IDENTITY; dfs_up(1, 0);  read ans_[v].
// NOTE the inner sibling loop is O(deg(v)); total O(sum deg^2). When Phi is
// INVERTIBLE replace it with subtract-the-child (see 15.1 solved). When it is
// NOT (max / signed product), use the prefix/suffix merge of 15.3 to get O(n).
\end{lstlisting}

\begin{complexitybox}
Two DFS passes over the tree, $\bigO{n}$ node-visits. The only per-node cost
beyond that is the sibling aggregate in \coderef{c20-twodfs}: as written it is
$\bigO{\deg(v)}$ per child, i.e.\ $\bigO{\deg(v)^2}$ per node and
$\bigO{\sum_v \deg(v)^2}$ overall -- up to $\bigO{n^2}$ for a star. Section
\secref{c20-prefixsuffix} buys back linearity; but whenever $\Phi$ is
\emph{invertible} the sibling aggregate is a single subtraction and each node
is $\bigO{\deg(v)}$, so the whole reroot is $\bigO{n}$. Memory $\bigO{n}$ for
the three arrays; stack depth $h$ -- at $n \ge 2\cdot10^{5}$ with chains legal,
run both passes on Ch.~12's explicit two-phase stack (\coderef{c09-iter}).
\end{complexitybox}

\begin{figure}[H]
\centering
\begin{tikzpicture}[diagram, scale=0.86, transform shape,
    tnode/.style={circle, draw=primaryblue, line width=0.9pt, fill=primaryblue!8,
                  minimum size=8mm, font=\small\bfseries, text=primaryblue!60!black}]
  \node[tnode] (n1) at (0,0)     {$1$};
  \node[tnode] (n2) at (-3,-1.6) {$2$};
  \node[tnode] (n3) at ( 3,-1.6) {$3$};
  \node[tnode] (n4) at (-5.2,-3.2) {$4$};
  \node[tnode] (n5) at (-1.6,-3.2) {$5$};
  \node[tnode] (n6) at ( 1.6,-3.2) {$6$};
  \node[tnode] (n7) at ( 5.2,-3.2) {$7$};
  \node[tnode] (n8) at (-1.6,-4.8) {$8$};
  \node[tnode] (n9) at ( 5.2,-4.8) {$9$};
  \node[tnode] (n10) at ( 5.2,-6.4) {$10$};
  \draw[tedge] (n1)--(n2) (n1)--(n3) (n2)--(n4) (n2)--(n5)
               (n3)--(n6) (n3)--(n7) (n5)--(n8) (n7)--(n9) (n9)--(n10);
  % size badge (down) top-left, F badge (answer) bottom-right
  \foreach \nd/\sz/\fv in {
      n1/10/30, n2/4/32, n3/5/30, n4/1/40, n5/2/38,
      n6/1/38, n7/3/34, n8/1/46, n9/2/40, n10/1/48} {
    \node[font=\scriptsize\ttfamily, text=successgreen!50!black, anchor=south east]
      at (\nd.north east) {$\sz$};
    \node[font=\tiny\ttfamily, text=warningred!75!black, anchor=north]
      at (\nd.south) {$F{=}\fv$};}
  \node[annot, font=\scriptsize, text=successgreen!50!black, anchor=south west]
    at (-0.35,0.5) {$\size(v)$ = down};
  \node[annot, font=\scriptsize, text=warningred!75!black, anchor=south]
    at (-6.4,-6.2) {$F(v)$ = answer for root $v$};
\end{tikzpicture}
\caption{The two passes on $T$ for CF~1187E. Green (top-right of each node) is
$\mathit{down}[v]=\size(v)$ from the post-order pass; red below is
$\mathit{dp}[v]=F(v)=\Phi(\mathit{down}[v],\mathit{up}[v])$, the sum of subtree
sizes when the tree is rooted at $v$, from the pre-order pass via the closed
edge-transfer $F(c)=F(v)+n-2\size(c)$. The maximum, $F(10)=48$, is CF~1187E's
answer on $T$ -- reproduced as \exref{c20-twodfs-run}.}
\label{fig:c20-updown}
\end{figure}

\subsection{Worked Example: CF 1187E's Fold on $T$}

\begin{example}{Sum-of-subtree-sizes rerooted across $T$}{c20-twodfs-run}
Reuse $T$ from \figref{c20-runtree-pair} ($n=10$, root $1$). The quantity is
$F(v) = \sum_{x} \size_{\text{rooted at }v}(x)$ -- the sum over all nodes of
their subtree size when $v$ is the root. Down pass: $\mathit{down}[v]=\size(v)$
in the reference rooting (Ch.~12's table). $F(1)=\sum_v\size(v)=30$. Up pass: the
closed edge-transfer $F(c)=F(v)+n-2\size(c)$ with $n=10$. Every value below is
hand-computed and reused, unchanged, by the trace box.
\end{example}

\begin{center}
\small
\renewcommand{\arraystretch}{1.2}
\begin{tabular}{l|cccccccccc}
\toprule
node $v$ & 1 & 2 & 3 & 4 & 5 & 6 & 7 & 8 & 9 & 10 \\
\midrule
$\size(v)=\mathit{down}[v]$ & 10 & 4 & 5 & 1 & 2 & 1 & 3 & 1 & 2 & 1 \\
$\mathit{up}[v]=F(v)-\mathit{down}[v]$ & 20 & 28 & 25 & 39 & 36 & 37 & 31 & 45 & 38 & 47 \\
$\mathit{dp}[v]=F(v)$ & 30 & 32 & 30 & 40 & 38 & 38 & 34 & 46 & 40 & \textbf{48} \\
\bottomrule
\end{tabular}
\end{center}

\noindent Sanity: $F(1)=30$ is the sum of the $\size$ row
($10{+}4{+}5{+}1{+}2{+}1{+}3{+}1{+}2{+}1$). The maximum is $F(10)=48$. Root the
tree at the leaf~$10$ and re-sum subtree sizes directly to check it without any
rerooting: the sizes become $10{+}9{+}8{+}7{+}5{+}4{+}1{+}2{+}1{+}1=48$
\checkmark\ -- an independent verification, which is the whole point of a
running example.

\begin{tracebox}[Edge-transfer on $T$, children visited in numeric order, $n=10$]
down pass (post-order 4,8,5,2,10,9,7,6,3,1): sizes are the table above;
$F(1)=\sum\size=30$.\\
up pass, $F(c)=F(v)+10-2\size(c)$:\\
$1\!\to\!2$: $30+10-2{\cdot}4=32$;\quad $1\!\to\!3$: $30+10-2{\cdot}5=30$.\\
$2\!\to\!4$: $32+10-2{\cdot}1=40$;\quad $2\!\to\!5$: $32+10-2{\cdot}2=38$.\\
$3\!\to\!6$: $30+10-2{\cdot}1=38$;\quad $3\!\to\!7$: $30+10-2{\cdot}3=34$.\\
$5\!\to\!8$: $38+10-2{\cdot}1=46$.\\
$7\!\to\!9$: $34+10-2{\cdot}2=40$.\quad $9\!\to\!10$: $40+10-2{\cdot}1=48$.\\
Answers $F=[30,32,30,40,38,38,34,46,40,48]$. $\max=48$ at $v=10$ \checkmark
(matches \figref{c20-updown} and the direct re-sum). Note every value is
derived from its \emph{parent's} already-known $F$, never recomputed from
scratch -- one pass, $\bigO{n}$. \checkmark
\end{tracebox}

\subsection{Variations and Twists}

\begin{itemize}[nosep]
  \item \textbf{Invertible $\Phi$ (sum, XOR):} recover "everything but child
        $c$" as $\mathit{down}[v] \ominus \mathit{down}[c]$ (subtract / XOR
        back), turning the sibling loop into $O(1)$ and the whole reroot into a
        clean $\bigO{n}$. This is what CF~1187E exploits. \hfill\lv{2}
  \item \textbf{Non-invertible $\Phi$ ($\max$, $\min$, signed product):} you
        cannot subtract; the sibling aggregate needs the prefix/suffix merge of
        \secref{c20-prefixsuffix}. \hfill\lv{4}
  \item \textbf{Richer per-node state (a pair or a count):} the template is
        unchanged, only the meaning of $\mathit{down}/\mathit{up}$ grows --
        \secref{c20-updown}. \hfill\lv{3}--\lv{4}
  \item \textbf{Reroot a game verdict:} $\Phi$ becomes a
        $\operatorname{mex}$/win-lose fold; the same two passes evaluate every
        root -- \secref{c20-counting} and the Summit. \hfill\lv{5}
\end{itemize}

\subsection{Common Pitfalls}

\begin{warningbox}[Trap 1 -- wrong up-pass seed]
$\mathit{up}[\text{root}]$ must be $\Phi$'s \emph{identity}, not $0$ by reflex.
For a sum, the identity is $0$ and you get lucky; for a $\max$-fold it is
$-\infty$, for a product it is $1$, for a count it is $0$. Seed a $\max$-reroot
with $0$ and every node above the true answer's argmax silently reports the
wrong value -- it will still pass hand-samples of all-positive trees.
\end{warningbox}

\begin{warningbox}[Trap 2 -- re-deriving $F(c)$ from the wrong $\size(c)$]
The closed form $F(c)=F(v)+n-2\size(c)$ uses $\size(c)$ measured in the
\emph{reference} rooting, i.e.\ the child's subtree size while $v$ is $c$'s
parent. Feed it the parent-side size and the sign of the correction flips; the
symptom is every non-root answer off by a constant that grows down the tree.
\end{warningbox}

\begin{warningbox}[Trap 3 -- recursion depth and the $F(1)$ accumulator]
At $n \le 2\cdot10^{5}$ with a chain legal, both DFS passes recurse $n$ deep
and overflow the default stack (a \re, not a \tle); run them iteratively with
\coderef{c09-iter}. Separately, $F$ is a sum of $n$ sizes, so it reaches
$\bigO{n^2}\sim 4\cdot10^{10}$ -- keep $\mathit{down}$, $\mathit{up}$, and $F$
in \lstinline{long long}; \lstinline{int} overflows immediately.
\end{warningbox}

\subsection{Problem Lab}

\solvedproblem{c20-twodfs-solved}{Tree Painting}{Codeforces 1187E}
  {https://codeforces.com/problemset/problem/1187/E}{CF 2100 · L4}
  {two-DFS rerooting, edge-transfer rule, sum of subtree sizes, invertible
   combine, long long}

\begin{problemstatement}[Tree Painting (statement and samples verified against the live problem)]
A tree of $n$ vertices ($2 \le n \le 2\cdot10^{5}$). All vertices start white.
On the first turn choose any one vertex and paint it black. On every later turn
choose a \emph{white} vertex adjacent to some black vertex and paint it black.
Each time you paint a vertex you gain points equal to the size of the connected
component of \emph{white} vertices that contains it at that moment. The game
ends when all vertices are black. Maximise the total points. \textbf{Example~1}:
$n=9$, edges $1\text{--}2,2\text{--}3,2\text{--}5,2\text{--}6,1\text{--}4,4\text{--}9,9\text{--}7,9\text{--}8$;
output \texttt{36}. \textbf{Example~2}: $n=5$, edges
$1\text{--}2,1\text{--}3,2\text{--}4,2\text{--}5$; output \texttt{14}.
\end{problemstatement}

\begin{attemptgate}[Attempt it yourself first -- \lv{4}, give yourself 30 minutes]
Two facts unlock it. (a)~Fix the first-painted vertex as a \emph{root}; then any
legal painting order paints a node only after its parent, so when you paint a
node its still-white component is exactly its \emph{subtree}. (b)~The total is
therefore $\sum_x \size(x)$ over the whole tree rooted at your first choice --
so maximise that sum over all roots, and \emph{moving the root one edge has a
closed-form effect on the sum}.
\peek{Rooted at $r$, what white component does a node $x$ get painted within,
      and how does that tie the score to subtree sizes?}
\nudge{Painting $x$ after its parent means the component still containing $x$
      is $\subtree(x)$: you gain $\size(x)$. Total $=F(r)=\sum_x\size(x)$.
      Now: when the root hops from $v$ to a child $c$, which nodes grow their
      subtree and which shrink?}
\fullhint{Re-rooting $v\to c$ keeps every node's subtree identical except that
$c$'s whole former subtree and the rest swap sizes: precisely $\size(c)$ nodes
lose $1$ from their subtree and the other $n-\size(c)$ gain $1$, so
$F(c)=F(v)+(n-\size(c))-\size(c)=F(v)+n-2\size(c)$. Compute $\size(\cdot)$ and
$F(1)=\sum\size$ in one post-order DFS; walk the tree top-down applying the
formula; take $\max_v F(v)$. Everything fits in \lstinline{long long}.}
\end{attemptgate}

\solstep{Understanding the Problem}
The scoring rule looks dynamic (it depends on \emph{when} you paint each node),
but the "adjacent to a black vertex" constraint forces the painted set to stay
connected, so the first vertex is a root and every node is painted after its
parent. Under that rooting, the white component around a node at the instant it
is painted is exactly its subtree -- so the points gained by that node are
$\size(x)$ \emph{regardless of the order} among valid orders. The whole game
therefore reduces to a single number per choice of first vertex:
$F(r)=\sum_x \size_{r}(x)$. Over $n$ roots that naively re-roots in
$\bigO{n^2}$; the edge-transfer of \thmref{c20-edge-transfer} collapses it to
$\bigO{n}$.

\solstep{Key Observations}
\begin{observation}
Fixing the first-painted vertex as the root, the total score is
$F(r)=\sum_{x}\size(x)$ -- the sum of subtree sizes in the tree rooted at $r$.
Paint order does not matter; only the root does.
\end{observation}
\begin{observation}
Re-rooting along edge $(v,c)$ where $c$ is $v$'s child in the reference rooting:
the two sides of the cut, $A=\subtree(c)$ (size $\size(c)$) and
$B=V\setminus\subtree(c)$ (size $n-\size(c)$), swap their contribution --
each node of $B$ gains one subtree member (the vertex $c$'s side reattaches),
each node of $A\setminus\{c\}$ loses the parent side. Net
$F(c)=F(v)+\card{B}-\card{A}=F(v)+n-2\size(c)$.
\end{observation}
\begin{observation}
$\Phi$ here is a plain sum, which is \emph{invertible}: the closed form above
is the "subtract the child" shortcut, so the sibling aggregate never needs a
loop and the whole reroot is $\bigO{n}$ -- no prefix/suffix required for this
problem.
\end{observation}

\solstep{Think It Through With Me}
\begin{thinkbox}
\deadend{Run Ch.~13's DFS once per vertex -- $n$ roots, each $\bigO{n}$ -- and
take the max. It is simple and correct.}
{Correct, and exactly $\bigO{n^2}$: at $n=2\cdot10^{5}$ that is $4\cdot10^{10}$
operations -- a certain \tle. The reroot is the \emph{whole} point of the
problem; brute force is the trap that looks like a solution.}
\deadend{Dynamically track the white components as nodes are painted, with a DSU.}
{The paint order is not part of the answer -- every valid order yields the same
$F(r)$ for a fixed first vertex $r$. Simulating the game re-introduces the
order you just proved irrelevant, and still leaves you choosing $r$.}
\aha{Reduce to $F(r)=\sum\size$; reroot with $F(c)=F(v)+n-2\size(c)$ in one
top-down pass after one bottom-down $\size$ pass. $\bigO{n}$, \lstinline{long
long}. That is the two-DFS template, first instantiation.}
\end{thinkbox}

\solstep{Solution}
Root at $1$. Post-order DFS: $\size[v]=1+\sum_{c}\size[c]$, and accumulate
$F(1)=\sum_v\size[v]$. Pre-order DFS: for each child $c$ of $v$,
$F(c)=F(v)+n-2\size[c]$, then recurse. Output $\max_{v} F(v)$. Two arrays, two
DFS passes, exactly \coderef{c20-twodfs} with $\Phi=+$ and
$\Phi\text{-shift}(\cdot)=+(\,n-2\size(c)\,)$.

\solstep{Proof of Correctness}
\begin{proofbox}[\coderef{c20-1187e-soln} returns $\max_r F(r)$]
\emph{Score $=F(r)$.} For a fixed first vertex $r$, colour $r$ as root. A legal
order paints each node after its parent, so at the moment node $x$ is painted
its white component is precisely $\subtree(x)$: it contains $x$ and all of
$x$'s unpainted descendants, and is separated from the rest by $x$'s already-
black parent. Hence $x$ contributes $\size(x)$ and the total is $F(r)$,
independent of order. \emph{Edge-transfer.} Re-rooting across $(v,c)$
partitions $V$ into $\subtree(c)$ (size $s=\size(c)$) and its complement
(size $n-s$). Every node $y$ with its subtree wholly inside one part keeps its
subtree size; the nodes whose subtree $=$ the whole of one part relative to the
other change by exactly $\pm1$: the $n-s$ nodes outside gain $+1$ and the $s$
nodes of $\subtree(c)$ lose $1$ (differentially, $c$ itself is neutral and the
correction telescopes to $n-2s$). So $F(c)=F(v)+n-2\size(c)$. \emph{Algorithm.}
The post-order pass computes all $\size[\cdot]$ and $F(1)$ by definition; the
pre-order pass applies the transfer parent-before-child, so by induction along
the pre-order every $F(v)$ is correct (\thmref{c20-edge-transfer}). Taking
$\max_v F(v)$ therefore returns $\max_r F(r)$. \qed
\end{proofbox}

\solstep{Implementation}
\begin{lstlisting}[caption={Codeforces 1187E -- the two-DFS template with Phi = + and shift = +(n - 2*size(c)); O(n), long long}, label={lst:c20-1187e-soln}]
#include <bits/stdc++.h>
using namespace std;
typedef long long ll;
int n;
vector<vector<int>> adj;
vector<ll> sz;     // subtree size: the down pass
vector<ll> F;      // sum of subtree sizes rooted at v: the answer per root

void dfs_down(int v, int p) {          // post-order: children before parent
    sz[v] = 1;
    for (int c : adj[v]) if (c != p) {
        dfs_down(c, v);
        sz[v] += sz[c];
    }
}
void dfs_up(int v, int p) {            // pre-order: parent before children
    for (int c : adj[v]) if (c != p) {
        F[c] = F[v] + (ll)n - 2 * sz[c];   // edge-transfer, thm:c20-edge-transfer
        dfs_up(c, v);
    }
}
int main() {
    scanf("%d", &n);
    adj.assign(n + 1, {});
    for (int i = 0; i < n - 1; i++) {
        int u, v; scanf("%d %d", &u, &v);
        adj[u].push_back(v); adj[v].push_back(u);
    }
    sz.assign(n + 1, 0); F.assign(n + 1, 0);
    dfs_down(1, 0);
    for (int v = 1; v <= n; v++) F[1] += sz[v];   // F(root) = sum of sizes
    dfs_up(1, 0);
    ll ans = 0;
    for (int v = 1; v <= n; v++) if (F[v] > ans) ans = F[v];
    printf("%lld\n", ans);
    return 0;
}
\end{lstlisting}

\solstep{Trace}
\begin{tracebox}[Official Example~2: $n=5$, edges $1\text{--}2,1\text{--}3,2\text{--}4,2\text{--}5$ $\Rightarrow$ 14]
Root $1$. Children: $1\!\to\!\{2,3\}$, $2\!\to\!\{4,5\}$. Down pass sizes:
$\size=[5,3,1,1,1]$ for $v=1..5$. $F(1)=5{+}3{+}1{+}1{+}1=11$.\\
Up pass, $F(c)=F(v)+5-2\size(c)$:\\
$1\!\to\!2$: $11+5-2{\cdot}3=10$;\quad $1\!\to\!3$: $11+5-2{\cdot}1=14$.\\
$2\!\to\!4$: $10+5-2{\cdot}1=13$;\quad $2\!\to\!5$: $10+5-2=13$.\\
$F=[11,10,14,13,13]$; $\max=14$ at node~$3$ \checkmark\ matches CF Example~2.
\end{tracebox}

\begin{tracebox}[Official Example~1 folds the same rule: $n=9$ $\Rightarrow$ 36]
Root $1$, $\size=[9,4,1,4,1,1,3,1,1]$ for $v=1..9$ (order $1,2,3,4,5,6,9,7,8$
by structure; $\size(2)=4$ sub-tree $\{2,3,5,6\}$, $\size(4)=4$ subtree
$\{4,9,7,8\}$, $\size(9)=3$). $F(1)=\sum\size=25$. Down the tree with
$F(c)=F(v)+9-2\size(c)$: $F(2)=26$, $F(4)=26$; $F(3)=F(5)=F(6)=33$;
$F(9)=F(4)+9-2{\cdot}3=29$; $F(7)=F(8)=F(9)+9-2{\cdot}1=36$. $\max=36$
at a leaf under~$9$ \checkmark\ matches CF Example~1.
\end{tracebox}

\solstep{Complexity}
\begin{complexitybox}
Time $\bigO{n}$: two DFS passes, and because $\Phi$ is a sum (invertible) the
per-node work is $O(\deg(v))$ with no sibling loop -- the closed transfer is
$\bigO{1}$ per edge. For $n=2\cdot10^{5}$ this is instantaneous. Memory
$\bigO{n}$ (adjacency $+\ \mathit{sz}+\ F$). The accumulator and both arrays
are \lstinline{long long}: $F$ peaks at $\bigO{n^2}\approx 4\cdot10^{10}$. Stack
depth is $h$; on a legal chain, convert the two passes to
\coderef{c09-iter}'s explicit stack or the run \re's. Verdict: \ac; the
archetypal two-DFS reroot -- down = subtree sizes, up = the $+n-2\size$ edge-
transfer, read $\max$ over all roots.
\end{complexitybox}

\solstep{Pitfalls and Variants}
\begin{warningbox}
Ways to lose CF~1187E: (i)~\lstinline{int} accumulators -- $F$ overflows past
$2\cdot10^{9}$; (ii)~using $n-2\size(c)$ with the \emph{parent-side} size (the
complement) instead of the child's -- flips the sign; (iii)~recursing the two
passes on a chain without raising the stack -- \re; (iv)~trying to simulate the
paint order with a DSU -- the order is irrelevant, and you will pay for the
dynamic structure. Variant: replace $F(v)$ with $\sum_v \depth(x)$
("sum of distances from $v$"); since $F(r)=n+\sum_x\depth(x)$ the very same
transfer applies -- a rerooting family (the practice below, CF~1881F, is one
distance-flavoured cousin).
\end{warningbox}

\subsection{Practice Problems}

\begin{stuckbox}
\begin{itemize}[nosep, leftmargin=1.4em]
  \item Did you reduce the game to $F(r)=\sum\size$ before touching a reroot --
        i.e.\ prove the paint order does not matter?
  \item Is your edge-transfer using $\size(c)$ from the \emph{reference} rooting
        (the child's own subtree), not the complement?
  \item Are $\mathit{sz}$ and $F$ both \lstinline{long long}, and are the two
        DFS passes safe against a depth-$n$ chain?
\end{itemize}
\end{stuckbox}

\begin{practiceset}
  \practice{Minimum Maximum Distance}{Codeforces 1881F}
    {https://codeforces.com/problemset/problem/1881/F}{CF 1700 · L3}
    {Given a tree and a set of $k$ marked vertices, for \emph{every} vertex $v$
    report $\max_{\text{marked } u} \dist(v,u)$, then take the minimum over $v$.
    A rerooting of the "farthest marked descendant / farthest via the parent"
    pair: $\mathit{down}[v]$ is the deepest mark inside $\subtree(v)$ and
    $\mathit{up}[v]$ the deepest mark reached through $v$'s parent; the
    edge-transfer is a $\max$ over $\mathit{down}$ of siblings plus
    $\mathit{up}[v]$, shifted by one edge -- Ch.~13's depth-histogram fold
    (\thmref{c10-depth-xor}) with \thmref{c20-edge-transfer} bolted on. Samples
    (verified): solve two test cases each; output one integer per case.
    }
\end{practiceset}

\paragraph{Back-cites (URLs live elsewhere in the book).}
\begin{itemize}[nosep,leftmargin=1.4em]
  \item \textbf{CF~1324F ``Maximum White Subtree''} (\chref{c10} owns this URL):
        the classic ``white minus black count for every root'' is
        \thmref{c20-edge-transfer} with $\Phi=+$ and the same $+n-2s$ shape;
        read Ch.~13's one-root version first, then reroot it here.
  \item \textbf{CF~342E / CF~219D} (\chref{c09} owns these URLs): distance-sum
        and "re-root a directed-edge count" exercises -- both are this template;
        their URLs live in Chapter~12, so cite by prose, not by link.
\end{itemize}

\begin{recallbox}[10-second refresher -- 15.1 The two-DFS template]
Split every node's rooted answer into $\mathit{down}[v]$ (Ch.~13's fold over
$\subtree(v)$) and $\mathit{up}[v]$ (the fold over everything above $v$). The
down pass is post-order; the up pass is pre-order and applies
\thmref{c20-edge-transfer}: $\mathit{up}[c]$ is the combine of $v$'s self,
$\mathit{up}[v]$, and $c$'s siblings' $\mathit{down}$, shifted per edge. Seed
$\mathit{up}[\text{root}]=\Phi\text{'s identity}$. CF~1187E
(\probref{c20-twodfs-solved}) uses $\Phi=+$ so the transfer is the closed
$F(c)=F(v)+n-2\size(c)$; on $T$ the answer is $\max F = 48$ at node~$10$.
\selftest{
  \item Why must the up pass run parent-before-child rather than post-order?
  \item For a $\max$-fold, what is the correct seed for $\mathit{up}[\text{root}]$?
  \item In CF~1187E, name the two nodes whose subtree sizes drive the
        $\pm(n-2\size(c))$ correction and say which side gains.
  \item What makes the sibling aggregate $\bigO{1}$ here but only $\bigO{\deg}$
        in \coderef{c20-twodfs} as written?
  \item State the leaf/base convention for the reference root's $\mathit{up}$.
}
\end{recallbox}
